\clearpage\begingroup
\SourceLeaf{178}\NoSourceHead
\begin{center}
\textbf{\Large [ELEKTRIZITÄT UND MAGNETISMUS.]}\par\bigskip
\textbf{\large\textls[120]{BRIEFWECHSEL.}}\par\bigskip
[1.]\par\medskip
\textbf{\textsc{Gauss} an \textsc{W. Weber}.}
\end{center}
\noindent\hspace*{\parindent}Wollen Sie, lieber \textsc{Weber}, einmal versuchen, ob Sie eine telegraphische Probe, die ich 5 Minuten nach dem Uhrzeichen anfangen will, lesen können? Hier der Schlüssel:
\[
\renewcommand{\arraystretch}{1.2}\setlength{\arraycolsep}{12pt}
\begin{array}{c||c|c|c|c|c|c}
\hline\hline
{}&-3&-2&-1&+1&+2&+3\\
\hline\hline
-3&y&=&w&f&=&=\\
-2&-&h&a&s&b&-\\
-1&k&d&e&i&g&p\\
+1&v&o&n&r&l&q\\
+2&-&c&t&u&m&-\\
+3&-&-&x&z&-&-
\end{array}
\]
\noindent Die offenen Stellen können für die Ziffern und etwa den Schlusspunkt reservirt bleiben. Es ist also z. B. $x=+3-1$, d. i. drei positive und nach einer kleinen Pause von etwa 4 Secunden Ein negativer Stoss. Positiv und negativ verstehe ich nach dem Magnetometer der Sternwarte; ob der Ihrige damit conform ist, oder ob Sie sämmtliche Zeichen entgegengesetzt rechnen müssen, erkennen Sie an dem heutigen Uhrzeichen, welches in der Sprache der Sternwarte $+$ sein soll.
\par\medskip
\noindent\hspace*{\parindent}[Göttingen] 16. Jul. 1835.\hfill\begin{tabular}[t]{@{}l@{}}Von Herzen\\\hspace*{4em}Ihr G[\textsc{auss}.]\end{tabular}\hspace*{\parindent}
\clearpage\SourceLeaf{179}\HistoricalHead{174}{L}{PHYSIK.\quad BRIEFWECHSEL.}
\noindent\hspace*{\parindent}Wenn Sie die Phrase nicht bis zu Ende gut lesen können, so wird glaube ich die Schuld nur an der kurzen Schwingungsdauer Ihrer Nadel liegen, die deshalb bald in bedeutende Bewegung kommen muss, wenn ich auch die Stösse so kurz, wie es die Einrichtung hier bis jetzt nur verstattet, gebe. Braucht man erst sehr schwere Nadeln, deren Schwingungsdauer etwa 1 Minute beträgt, so wird, hoffe ich, diese Methode nichts zu wünschen übrig lassen, und [es werden] in Einer Minute füglich 4 Buchstaben übermacht werden können.
\clearpage\endgroup
